\documentclass{beamer}

% Theme selection
\usetheme{Madrid}
\usecolortheme{default}

% Metadata
\title{Modeling Aspects of AI}
\subtitle{Foundations, Architectures, and Ethics}
\author{Prepared by Gemini}
\date{\today}

\begin{document}

% Title Page
\begin{frame}
    \titlepage
\end{frame}

% Table of Contents
\begin{frame}{Overview}
    \tableofcontents
\end{frame}

% Section 1: Data Modeling
\section{Data Representation}
\begin{frame}{Data Representation}
    \begin{itemize}
        \item \textbf{Vector Space Model:} Representing tokens in high-dimensional space.
        \item \textbf{Embeddings:} Capturing semantic meaning through proximity.
        \item \textbf{Normalization:} Ensuring data consistency for gradient descent.
    \end{itemize}
    \begin{block}{The Objective Function}
        Typically, we aim to minimize a loss function $L(\theta)$ such as:
        $$L(\theta) = \frac{1}{n} \sum_{i=1}^{n} (y_i - f(x_i; \theta))^2$$
    \end{block}
\end{frame}

% Section 2: Neural Architectures
\section{Neural Architectures}
\begin{frame}{Key Modeling Architectures}
    \begin{columns}
        \column{0.5\textwidth}
        \textbf{Transformers}
        \begin{itemize}
            \item Self-Attention mechanism
            \item Parallel processing
            \item Scalability
        \end{itemize}
        
        \column{0.5\textwidth}
        \textbf{Diffusion Models}
        \begin{itemize}
            \item Forward noise injection
            \item Reverse denoising process
            \item Generative prowess
        \end{itemize}
    \end{columns}
\end{frame}

% Section 3: Optimization
\section{Optimization Aspects}
\begin{frame}{Optimization and Convergence}
    The goal of AI modeling is finding the global minimum of the error surface.
    \begin{itemize}
        \item \textbf{Stochastic Gradient Descent (SGD):} $ \theta_{t+1} = \theta_t - \eta \nabla L(\theta_t) $
        \item \textbf{Regularization:} Preventing overfitting via $L1$ or $L2$ penalties.
        \item \textbf{Hyperparameter Tuning:} Learning rates, batch sizes, and dropout.
    \end{itemize}
\end{frame}

% Conclusion
\section{Conclusion}
\begin{frame}{Summary}
    \begin{itemize}
        \item Modeling requires a balance of \textbf{architecture}, \textbf{data}, and \textbf{compute}.
        \item Future trends: Neuro-symbolic AI and energy-efficient modeling.
    \end{itemize}
    \vspace{1cm}
    \centering
    \Large \textbf{Questions?}
\end{frame}

\end{document}
% Go to Overleaf.com, create a new Blank Project, and paste the code above. Hit Recompile, and your PDF will appear on the right.